% BEGIN VOL07-EXPANSION VII35
\section{Supplementary worked examples}
\label{sec:vii35-pedagogy}

\begin{example}[Modular cusp]\label{ex:vii35-ped-01}
For \(PSL(2,\mathbb Z)\), the translation \(T:z\mapsto z+1\) stabilizes \(\infty\).  The region \(0\le x\le1\), \(y\ge Y\) quotients to a cusp cylinder whose hyperbolic area is \(1/Y\).
\end{example}

\begin{example}[A hyperbolic matrix]\label{ex:vii35-ped-02}
The matrix \(A=\begin{pmatrix}2&1\\1&1\end{pmatrix}\) lies in \(SL(2,\mathbb Z)\) and has trace \(3\).  It is hyperbolic and determines a quotient closed geodesic of primitive length \(2\operatorname{arcosh}(3/2)\) when primitive in the group.
\end{example}

\begin{example}[Fundamental-domain side pairing]\label{ex:vii35-ped-03}
In the modular fundamental domain, \(T\) pairs the vertical sides and \(S:z\mapsto-1/z\) pairs the circular arcs.  Group generators are therefore visible directly as boundary identifications of a geometric tile.
\end{example}

\section{Graded supplementary exercises}
The following exercises add a deliberate progression from concrete calculation to structural proof, hypothesis testing, application, and synthesis.

\subsection*{Standard computations and constructions}

\begin{exercise}[Discrete definition]\label{exr:vii35-ped-01}
What two ingredients define a Fuchsian group?
\end{exercise}

\begin{hint}
Name ambient group and topology.
\end{hint}

\begin{solution}
It is a discrete subgroup of \(PSL(2,\mathbb R)\).
\end{solution}

\begin{exercise}[Modular generators]\label{exr:vii35-ped-02}
Name standard generators of \(PSL(2,\mathbb Z)\).
\end{exercise}

\begin{hint}
Use translation and inversion.
\end{hint}

\begin{solution}
\(T(z)=z+1\) and \(S(z)=-1/z\).
\end{solution}

\begin{exercise}[Matrix type]\label{exr:vii35-ped-03}
Classify \(\begin{pmatrix}2&1\\1&1\end{pmatrix}\).
\end{exercise}

\begin{hint}
Its trace is \(3\).
\end{hint}

\begin{solution}
Hyperbolic.
\end{solution}

\begin{exercise}[Closed length]\label{exr:vii35-ped-04}
Compute its translation length symbolically.
\end{exercise}

\begin{hint}
Use the trace-length relation.
\end{hint}

\begin{solution}
\(2\operatorname{arcosh}(3/2)\).
\end{solution}

\begin{exercise}[Cusp area]\label{exr:vii35-ped-05}
Compute the area of \([0,1]\times[Y,\infty)\) in the half-plane metric.
\end{exercise}

\begin{hint}
Integrate \(y^{-2}\).
\end{hint}

\begin{solution}
\(1/Y\).
\end{solution}

\subsection*{Proofs}

\begin{exercise}[Proper discontinuity]\label{exr:vii35-ped-06}
Why does discreteness imply only finitely many translates of a compact set meet it?
\end{exercise}

\begin{hint}
Otherwise produce group elements accumulating near an isometry moving points only slightly.
\end{hint}

\begin{solution}
Compactness plus infinitely many intersections would yield an accumulating sequence in the isometry group, contradicting discreteness.
\end{solution}

\begin{exercise}[Torsion-free quotient]\label{exr:vii35-ped-07}
Why is a torsion-free Fuchsian quotient a smooth surface rather than an orbifold?
\end{exercise}

\begin{hint}
Interior stabilizers correspond to elliptic torsion.
\end{hint}

\begin{solution}
Torsion-free discreteness makes the action free and properly discontinuous, so the projection is a covering and local isometry.
\end{solution}

\begin{exercise}[Axis projection]\label{exr:vii35-ped-08}
Why does a hyperbolic element's axis project to a closed geodesic?
\end{exercise}

\begin{hint}
The element translates the axis by one period.
\end{hint}

\begin{solution}
Quotienting the invariant axis by the cyclic translation identifies points separated by the translation length, producing a closed geodesic.
\end{solution}

\begin{exercise}[Cocompact cusp]\label{exr:vii35-ped-09}
Why can a cocompact Fuchsian group have no cusp?
\end{exercise}

\begin{hint}
A cusp is a noncompact end.
\end{hint}

\begin{solution}
A compact quotient has no ends, hence no parabolic cusp neighborhoods.
\end{solution}

\subsection*{Counterexamples and hypothesis tests}

\begin{exercise}[Discrete means countable]\label{exr:vii35-ped-10}
Test the converse: every countable subgroup of \(PSL(2,\mathbb R)\) is discrete.
\end{exercise}

\begin{hint}
Countable sets can accumulate.
\end{hint}

\begin{solution}
False.  A countable subgroup may contain elements approaching the identity.
\end{solution}

\begin{exercise}[Elliptic allowed torsion-free]\label{exr:vii35-ped-11}
Test the claim.
\end{exercise}

\begin{hint}
Interior elliptic elements have finite-order stabilizer in a discrete Fuchsian group.
\end{hint}

\begin{solution}
False.  Torsion-free Fuchsian groups contain no nontrivial elliptic elements.
\end{solution}

\begin{exercise}[Finite area means compact]\label{exr:vii35-ped-12}
Test the claim using the modular surface.
\end{exercise}

\begin{hint}
Cusps may have finite area.
\end{hint}

\begin{solution}
False.  The modular quotient has finite area but is noncompact because of its cusp.
\end{solution}

\subsection*{Applications and investigations}

\begin{exercise}[Mesh-like quotient]\label{exr:vii35-ped-13}
Why are fundamental domains useful computationally for quotient geometry?
\end{exercise}

\begin{hint}
They reduce an infinite cover to one tile plus side-pairing rules.
\end{hint}

\begin{solution}
Algorithms can work on a finite region while applying group transformations whenever trajectories cross paired sides.
\end{solution}

\begin{exercise}[Length spectrum]\label{exr:vii35-ped-14}
How do hyperbolic conjugacy classes contribute to the length spectrum of a quotient surface?
\end{exercise}

\begin{hint}
Associate each primitive class with its axis.
\end{hint}

\begin{solution}
Each primitive hyperbolic conjugacy class determines a closed geodesic whose length equals the element's translation length.
\end{solution}

\subsection*{Challenge problems}

\begin{exercise}[Modular area]\label{exr:vii35-ped-15}
State the hyperbolic area of the standard modular fundamental domain.
\end{exercise}

\begin{hint}
Use the classical value.
\end{hint}

\begin{solution}
\(\pi/3\).
\end{solution}

\begin{exercise}[Cusp shrink]\label{exr:vii35-ped-16}
As \(Y\to\infty\), what happens to the area of the standard cusp tail above height \(Y\)?
\end{exercise}

\begin{hint}
Use \(1/Y\).
\end{hint}

\begin{solution}
It tends to zero even though the tail remains infinitely long in the vertical direction.
\end{solution}

% END VOL07-EXPANSION VII35
